\section{Proof systems}\label{sec:prelims-nizks}

Zero knowledge proof systems enable a prover to convince verifiers that a
statement is true without revealing any information beyond the validity of the
statement.

\paragraph{Syntax.} A non interactive zero knowledge proof (NIZK) system
consists of the following algorithms:

\begin{description}

  \item[$\crs \sample \NIZKgen(\secparam)$ :] The setup algorithm outputs a
    common reference string $\crs$.

  \item[$\proof \leftarrow \NIZKprove(\crs, \inst, \witness)$ :] The proving
    algorithm takes as input a common reference string $\crs$, statement
    $\inst$, and witness $\witness$, and outputs a proof $\proof$.

  \item[$\{0, 1\} \leftarrow \NIZKverify(\crs, \inst, \proof)$ :] The
    verification algorithm takes as input a common reference string $\crs$,
    statement $\inst$, and proof $\proof$, and outputs $0$ or $1$.

\end{description}

\paragraph{Correctness.} A non interactive zero knowledge proof system must
satisfy the following correctness requirement: for every common reference
string $\crs$ generated by $\NIZKgen(\secparam)$ and every statement and
witness pair $(\inst, \ \witness) \in \lang$, the following holds:
%
\begin{equation*}
%
  \NIZKverify(\crs, \inst, \NIZKprove(\crs, \inst, \witness)) = 1.
%
\end{equation*}
%
\paragraph{Security.} The constructions presented in this thesis require
different security notions for non interactive zero knowledge proof systems,
including zero knowledge, simulation soundness, and simulation extractability.
We define below the notions used throughout the thesis, and introduce any
construction-specific requirements in the following chapters as needed.

\paragraph{Labelled proofs.} Certain constructions require proofs that are bound
to a message to prevent proof replaying.  We can attach a label to each
proof to prevent this.  A label is an arbitrary string that is fixed at proving
time, supplied again at verification, and is not otherwise interpreted by the
proof system.  A labelled NIZK extends the syntax as follows:

\begin{description}

  \item[$\proof \leftarrow \NIZKprove(\crs, \inst, \witness, \lbl)$ :] The
    proving algorithm additionally takes a label $\lbl \in \{0,1\}^*$.

  \item[$\{0, 1\} \leftarrow \NIZKverify(\crs, \inst, \proof, \lbl)$ :] The
    verification algorithm additionally takes the label $\lbl$, and accepts
    only proofs produced under that label.

\end{description}

Correctness is as before but over all labels.  Labelled proofs are equivalent
to signatures of knowledge~\cite{C:ChaLys06}.

\paragraph{Simulation and extraction.} To state the security notions we
require, we introduce the following additional algorithms.

\begin{description}

  \item[$(\crs, \td) \sample \NIZKsimgen(\secparam)$ :] The simulated setup
    algorithm outputs a common reference string $\crs$ together with a trapdoor
    $\td$.  The distribution of $\crs$ is computationally indistinguishable
    from that produced by $\NIZKgen(\secparam)$.

  \item[$\proof \leftarrow \NIZKsimprove(\crs, \td, \inst, \lbl)$ :] The
    simulator takes the trapdoor, a statement $\inst$, and a label $\lbl$, and
    outputs a proof, without access to a witness.

  \item[$\witness \leftarrow \NIZKextract(\crs, \td, \inst, \proof, \lbl)$ :]
    The extractor takes the trapdoor, a statement, a proof, and a label, and
    outputs a candidate witness.

\end{description}

\paragraph{Zero knowledge.} A labelled NIZK is zero knowledge if simulated
proofs are indistinguishable from honestly generated ones.  We define the
following experiment between a challenger and an adversary $\adv$.

\begin{experimentbox}{Zero knowledge}{zk}

  \begin{enumerate}

    \item The challenger samples $b \sample \{0,1\}$.  If $b = 0$ it generates
      $\crs \sample \NIZKgen(\secparam)$. If $b = 1$ it generates $(\crs, \td)
      \sample \NIZKsimgen(\secparam)$.  It sends $\crs$ to $\adv$.

    \item $\adv$ is given oracle access to $\simoracle$, which on input
      $(\inst, \witness, \lbl)$ with $\rel(\inst, \witness) = 1$ returns
      $\NIZKprove(\crs, \inst, \witness, \lbl)$ if $b = 0$, and
      $\NIZKsimprove(\crs, \td, \inst, \lbl)$ if $b = 1$.

    \item $\adv$ outputs a bit $b'$.

  \end{enumerate}

  $\adv$ wins the experiment if $b' = b$.

\end{experimentbox}

The advantage of $\adv$ is defined as follows:
%
\begin{equation*}
%
  \Adv^{\mathsf{zk}}_{\mathsf{NIZK}, \adv}(\secparam) = \left| \Pr[b' =
  b] - \frac{1}{2} \right|.
%
\end{equation*}
%
We say the proof system satisfies the property of zero knowledge if this
advantage is small for every admissible adversary $\adv$.

\paragraph{Simulation extractability.} We require more than soundness: our
security arguments need not only that accepting proofs are for true statements,
but that a witness can be recovered from them, and that this remains so after
the adversary has observed simulated proofs.  Simulation extractability
captures both of these properties.

\begin{experimentbox}{Simulation extractability}{simext}

  \begin{enumerate}

    \item The challenger generates $(\crs, \td) \sample \NIZKsimgen(\secparam)$
      and sends $\crs$ to $\adv$.

    \item $\adv$ is given oracle access to $\simoracle$, which on input
      $(\inst, \lbl)$ returns $\proof \leftarrow \NIZKsimprove(\crs, \td,
      \inst, \lbl)$ and stores $(\inst, \lbl)$ in the set $\Query$.

    \item $\adv$ outputs $(\inst^*, \proof^*, \lbl^*)$.

    \item Challenger extracts $\witness^* \leftarrow \NIZKextract(\crs,
      \td, \inst^*, \proof^*, \lbl^*)$.

  \end{enumerate}

  $\adv$ wins the experiment if $\NIZKverify(\crs, \inst^*, \proof^*, \lbl^*) =
  1$, $(\inst^*, \lbl^*) \notin \Query$, and $\rel(\inst^*, \witness^*) \neq
  1$.

\end{experimentbox}

The simulation extractability advantage of $\adv$ is defined as follows:
%
\begin{equation*}
%
  \Adv^{\mathsf{simext}}_{\mathsf{NIZK}, \adv}(\secparam) = \Pr[\adv\text{
  wins}].
%
\end{equation*}
%
A labelled NIZK is simulation extractable if this advantage is small for
every admissible adversary $\adv$.

The condition $(\inst^*, \lbl^*) \notin \Query$ makes proofs non-malleable
across labels.

Zero knowledge and simulation extractability are independent properties.
Zero knowledge ensures that simulated proofs are indistinguishable from
honestly generated ones. Simulation extractability ensures that a witness can
be extracted from any verifying proof, even after the adversary has observed
simulated proofs. Since every proof system constructed in this thesis satisfies
both properties, we often refer to them simply as proofs.
